\documentclass{amsart}

\usepackage{amssymb}
\usepackage{mathtools}
\usepackage{booktabs}
\usepackage{needspace}
\usepackage[shortlabels]{enumitem}
\usepackage[hidelinks]{hyperref}

\hypersetup{
  pdftitle={Type and relations of four-generated numerical semigroups},
  pdfauthor={Knowledge Sense},
  pdfsubject={Mathematics},
  pdfkeywords={numerical semigroup, embedding dimension four, type, minimal presentation, Betti number, Apery set, Scarf complex}
}

\newtheorem{theorem}{Theorem}[section]
\newtheorem{proposition}[theorem]{Proposition}
\newtheorem{lemma}[theorem]{Lemma}
\newtheorem{corollary}[theorem]{Corollary}
\theoremstyle{definition}
\newtheorem{definition}[theorem]{Definition}
\newtheorem{example}[theorem]{Example}
\theoremstyle{remark}
\newtheorem{remark}[theorem]{Remark}

\numberwithin{equation}{section}
\setlength{\emergencystretch}{2em}

\newcommand{\N}{\mathbb{N}}
\newcommand{\Z}{\mathbb{Z}}
\newcommand{\R}{\mathbb{R}}
\newcommand{\kk}{\Bbbk}
\newcommand{\fm}{\mathfrak{m}}
\newcommand{\fn}{\mathfrak{n}}
\newcommand{\cR}{\mathcal{R}}
\newcommand{\Ap}{\operatorname{Ap}}
\DeclareMathOperator{\PF}{PF}
\DeclareMathOperator{\Betti}{Betti}
\DeclareMathOperator{\Soc}{Soc}
\DeclareMathOperator{\supp}{supp}
\DeclareMathOperator{\Std}{Std}
\DeclareMathOperator{\Tor}{Tor}
\DeclareMathOperator{\lcm}{lcm}
\newcommand{\wt}{\ell}
\newcommand{\Apd}{\Ap(S,d_0)}

\begin{document}

\title{Type and relations of four-generated numerical semigroups}

\author{Knowledge Sense}
\address{Knowledge Sense}
\email{math@knowledgesense.jp}

\date{12 October 2026}

\subjclass[2020]{Primary 20M14; Secondary 13D02, 13F65, 05E40}

\keywords{Numerical semigroup, embedding dimension four, type, minimal presentation, Betti number, Ap\'ery set, Scarf complex}

\begin{abstract}
Let $S$ be a numerical semigroup with embedding dimension four, let $t(S)$ be its type and let $\eta(S)$ be the cardinality of a minimal presentation of $S$. We prove that
\[
t(S)+1\le \eta(S)\le 2t(S)+3 .
\]
The lower inequality confirms a conjecture of Chomicz, and the upper inequality improves his bound $\eta(S)\le 4t(S)+5$ and the earlier bound $\eta(S)\le 9t(S)+4$ of Bresinsky. Both inequalities are attained. The proof starts from Chomicz's description of the Ap\'ery set as a staircase of lattice points. We label the generators so that the defining ideal of the Artinian reduction either equals the monomial ideal of the staircase or differs from it by a single binomial. We compare $\eta(S)$ with the number $\mu$ of minimal generators of this monomial ideal, by bounding the number of these generators that become redundant, and $t(S)$ with the number $t'$ of corners of the staircase, through its top slice. Finally, we compare $\mu$ with $t'$ by counting the triangles of a Scarf triangulation that is sensitive to the supports of the generators.
\end{abstract}

\maketitle

\section{Introduction}

Let $\N$ denote the set of nonnegative integers. A \emph{numerical semigroup} is a submonoid $S\subseteq\N$ with finite complement. It has a unique minimal generating set, whose cardinality is the \emph{embedding dimension} $e(S)$. An integer $f\notin S$ is a \emph{pseudo-Frobenius number} of $S$ if $f+s\in S$ for every nonzero $s\in S$; the set of these numbers is $\PF(S)$, and $t(S)=|\PF(S)|$ is the \emph{type} of $S$. If $S$ is minimally generated by $d_0,\dots,d_n$, the \emph{factorization map} $\varphi\colon\N^{n+1}\to S$ is $\varphi(c)=c_0d_0+\dots+c_nd_n$, and a \emph{minimal presentation} of $S$ is a minimal generating set of the congruence $\{(c,c')\in\N^{n+1}\times\N^{n+1} : \varphi(c)=\varphi(c')\}$. Its cardinality depends only on $S$ and is denoted $\eta(S)$ \cite{Ros96,RGS09}. For a field $\kk$, the semigroup ring $\kk[S]=\kk[u^s : s\in S]\subseteq\kk[u]$ is a one-dimensional Cohen--Macaulay domain. If $\kk[S]$ is presented as a quotient of a polynomial ring in $n+1$ variables, then $\eta(S)$ is the minimal number of generators of the defining ideal \cite{Her70}, and $t(S)$ is the last Betti number of $\kk[S]$, that is, its Cohen--Macaulay type; see \cite{RGS09,Sta18}.

If $e(S)=2$, then $t(S)=\eta(S)=1$. If $e(S)=3$, Herzog \cite{Her70} showed that $\eta(S)=2$ and $t(S)=1$ when $S$ is symmetric, and $\eta(S)=3$ and $t(S)=2$ otherwise; in both cases $\eta(S)=t(S)+1$. In embedding dimension four both invariants are unbounded: Bresinsky \cite{Bre75a} found semigroups with arbitrarily large $\eta(S)$, and Backelin found semigroups with arbitrarily large $t(S)$ \cite[p.~75]{FGH87}. Bresinsky proved that $\eta(S)\le 9t(S)+4$ whenever $e(S)=4$ \cite[Theorem~7]{Bre88}. Moscariello and Sammartano asked whether, conversely, $t(S)$ is bounded by a function of $\eta(S)$ when $e(S)=4$ \cite[Problem~25]{MS25}. Chomicz answered this question in \cite[Theorem~1.4]{Cho26}, proving
\[
t(S)-11\ \le\ \eta(S)\ \le\ 4t(S)+5 \qquad (e(S)=4),
\]
and conjectured that $\eta(S)\ge t(S)+1$ \cite[Conjecture~5.1]{Cho26}. He also observed that this inequality fails in embedding dimension five \cite[Remark~5.2]{Cho26}. Bounds for the Betti numbers of $\kk[S]$ in terms of other invariants, such as the width, that is, the difference between the largest and the smallest minimal generator, are studied in \cite{CMS24}.

Our main result proves Chomicz's conjecture and improves the upper bound.

\begin{theorem}\label{thm:main}
Let $S$ be a numerical semigroup with embedding dimension four. Then
\[
t(S)+1\ \le\ \eta(S)\ \le\ 2t(S)+3 .
\]
\end{theorem}

Since $\kk[S]$ has projective dimension three over a polynomial ring in four variables, Theorem~\ref{thm:main} controls all its Betti numbers.

\begin{corollary}\label{cor:betti}
Let $S$ be a numerical semigroup with embedding dimension four, let $\kk$ be a field, and let $\beta_0,\dots,\beta_3$ be the total Betti numbers of $\kk[S]$ over a polynomial ring in four variables. Then $\beta_0=1$, $\beta_1=\eta(S)$, $\beta_3=t(S)$ and $\beta_2=\eta(S)+t(S)-1$. Consequently
\[
t(S)+1\le\beta_1\le 2t(S)+3,\qquad 2t(S)\le\beta_2\le 3t(S)+2,\qquad \tfrac12(\beta_1-3)\le\beta_3\le\beta_1-1 .
\]
\end{corollary}

In particular each of $\beta_1,\beta_2,\beta_3$ is bounded above and below by linear functions of each of the other two; compare \cite[Problem~24]{MS25} and the discussion in \cite[Section~1]{Cho26}.

\medskip\noindent\textbf{Sharpness.} Both inequalities of Theorem~\ref{thm:main} are attained. Chomicz lists the semigroups $\langle 77,85,89,131\rangle$, $\langle 76,99,143,223\rangle$ and $\langle 195,200,209,226\rangle$, for which $(t(S),\eta(S))$ equals $(3,4)$, $(6,7)$ and $(8,9)$, respectively \cite[Section~5]{Cho26}. For the upper inequality, recall that $t(S)=1$ if and only if $S$ is symmetric \cite{Kun70}, and that a symmetric semigroup with $e(S)=4$ satisfies $\eta(S)\in\{3,5\}$ by a theorem of Bresinsky \cite{Bre75b}; see also \cite[Theorem~4.18]{Cho26}. Hence every symmetric four-generated semigroup that is not a complete intersection has $\eta(S)=5=2t(S)+3$; an example is $S=\langle 5,6,7,8\rangle$, discussed in Example~\ref{ex:5678}. Thus the additive constant $3$ cannot be decreased as long as the coefficient $2$ is kept. We do not know whether the coefficient $2$ is optimal for large type. In the computations described in Section~\ref{sec:examples}, the equality $\eta(S)=2t(S)+3$ occurs only for $t(S)=1$.

\medskip\noindent\textbf{Strategy.} Write $S=\langle d_0,d_1,d_2,d_3\rangle$. We work with the Artinian reduction $A=\kk[S]/(u^{d_0})$, presented as $A=\kk[x,y,z]/J$. Its minimal number of relations is $\eta(S)$ and its socle dimension is $t(S)$. Chomicz \cite{Cho26} introduced a labelling of the generators, distinguishing \emph{primary} and \emph{secondary} semigroups, for which the Ap\'ery set $\Apd$ is represented by a staircase $T\subset\N^3$, and he described the minimal generators of the monomial ideal $I$ of this staircase. With this labelling, $J$ either equals $I$ or is obtained from $I$ by replacing one pure power $z^a$ by a binomial $z^a-p$ (Lemma~\ref{lem:J}). We compare $\eta(S)$ and $t(S)$ with the number $\mu$ of minimal generators of $I$ and the number $t'$ of corners of $T$. Let $r$ be the number of minimal generators of $I$ divisible by $xyz$; Chomicz's description gives $r\le1$. The proof has four ingredients.
\begin{enumerate}[(i)]
\item For every Artinian monomial ideal in three variables, $t'\le\mu+r-2$ (Proposition~\ref{prop:scarf}). This is a consequence of the generic deformations and the Scarf triangulations of Bayer, Peeva and Sturmfels \cite{BPS98}, applied to a deformation that keeps zero exponents zero.
\item The distinguished generator $d_0$ of a primary semigroup can be chosen so that $r=0$ or $J=I$ (Lemma~\ref{lem:choice}).
\item If $J\neq I$ and $r=0$, at most one minimal generator of $I$ is redundant among the generators of $J$ given by Lemma~\ref{lem:J}, so that $\eta(S)\ge\mu-1$ (Propositions~\ref{prop:primary} and~\ref{prop:secondary}).
\item If $J\ne I$, then $t'-t(S)\le r+1$ (Proposition~\ref{prop:topslice}).
\end{enumerate}
The lower bound follows from (i)--(iii). The upper bound follows from (ii), (iv) and Dung's inequality $\mu\le 2t'+1$ for Artinian monomial ideals in three variables \cite[Theorem~3.2]{Dun18}.

\medskip\noindent\textbf{Organization.} Section~\ref{sec:prelim} collects facts on factorizations, the Artinian reduction and monomial ideals in three variables, and proves Proposition~\ref{prop:scarf}. Section~\ref{sec:staircase} recalls Chomicz's staircase and its boundary and compares $J$ with $I$. Section~\ref{sec:choice} explains the choice of $d_0$. Sections~\ref{sec:relations} and~\ref{sec:type} compare $\eta(S)$ with $\mu$ and $t(S)$ with $t'$. Section~\ref{sec:proof} proves Theorem~\ref{thm:main} and Corollary~\ref{cor:betti}, and Section~\ref{sec:examples} contains examples and a description of supplementary computations.

\section{Preliminaries}\label{sec:prelim}

Throughout, $S$ is a numerical semigroup minimally generated by four positive integers $d_0,d_1,d_2,d_3$, and $\kk$ is an arbitrary field. The labelling of the generators will be fixed in Section~\ref{sec:staircase}; until then it is arbitrary.

\subsection{Factorizations and Betti elements}

For $s\in S$, the set $\varphi^{-1}(s)\subseteq\N^4$ is the set of \emph{factorizations} of $s$, and the \emph{support} of $c\in\N^4$ is $\supp(c)=\{i : c_i>0\}$. Two factorizations $c,c'$ of $s$ are \emph{$\cR$-related} if there is a sequence $c=c^1,\dots,c^q=c'$ in $\varphi^{-1}(s)$ such that $\supp(c^j)\cap\supp(c^{j+1})\neq\emptyset$ for all $j$. Let $\kappa(s)$ be the number of $\cR$-classes of $\varphi^{-1}(s)$. An element $s\in S$ is a \emph{Betti element} if $\kappa(s)\ge2$; the set of Betti elements is denoted $\Betti(S)$.

Let $Q=\kk[x_0,x,y,z]$ be graded by $\deg x_0=d_0$, $\deg x=d_1$, $\deg y=d_2$ and $\deg z=d_3$, let $\fn$ be its homogeneous maximal ideal, and let $I_S$ be the kernel of the surjection $Q\to\kk[S]$ sending $x_0,x,y,z$ to $u^{d_0},u^{d_1},u^{d_2},u^{d_3}$. The ideal $I_S$ is homogeneous and is spanned by binomials $x^c-x^{c'}$ with $\varphi(c)=\varphi(c')$, where $x^c=x_0^{c_0}x^{c_1}y^{c_2}z^{c_3}$.

\begin{proposition}[\cite{Her70,Ros96}]\label{prop:herzog-rosales}
For every $s\in S$, $\dim_\kk(I_S/\fn I_S)_s=\kappa(s)-1$. In particular $\eta(S)$ equals the minimal number of generators of $I_S$, and $(I_S/\fn I_S)_s\neq0$ if and only if $s\in\Betti(S)$.
\end{proposition}

Herzog \cite{Her70} showed that minimal presentations of $S$ correspond to minimal binomial generating sets of $I_S$, and Rosales \cite{Ros96} showed that every minimal presentation contains exactly $\kappa(s)-1$ pairs of degree $s$ for each $s$; see \cite[Theorem~3.1]{Cho26}, \cite{RGS09} and \cite{Sta18}.

\subsection{The Artinian reduction}

Let $P=\kk[x,y,z]$ with the grading induced from $Q$, and let $\fm=(x,y,z)$. For $v=(v_1,v_2,v_3)\in\N^3$ we write $x^v=x^{v_1}y^{v_2}z^{v_3}$ and
\[
\wt(v)=v_1d_1+v_2d_2+v_3d_3 ,
\]
so that $\wt(v)$ is the degree of $x^v$. We call $\wt(v)$ the \emph{weight} of $v$ and of $x^v$. Let $\varepsilon_1,\varepsilon_2,\varepsilon_3$ be the unit vectors of $\N^3$, and order $\N^3$ componentwise.

The \emph{Ap\'ery set} of $S$ with respect to $d_0$ is $\Apd=\{s\in S : s-d_0\notin S\}$. It contains exactly one element of each residue class modulo $d_0$, so $|\Apd|=d_0$ \cite{RGS09}. If $w\in\Apd$ and $w=w'+w''$ with $w',w''\in S$, then $w'\in\Apd$. We use the description \cite[Proposition~2.20]{RGS09} (see also \cite[(4)]{Cho26})
\begin{equation}\label{eq:PF}
\PF(S)=\{\,w-d_0 : w\in\Apd,\ w+d_i\notin\Apd \text{ for } i=0,1,2,3\,\}.
\end{equation}
The condition for $i=0$ is automatic.

Let $A=\kk[S]/(u^{d_0})$, and let $J\subseteq P$ be the kernel of the surjection $P\to A$ sending $x,y,z$ to the classes of $u^{d_1},u^{d_2},u^{d_3}$.

\begin{lemma}\label{lem:reduction}
The ideal $J$ and the ring $A$ have the following properties.
\begin{enumerate}[(a)]
\item For all $i\ge0$ and $s\in\Z$, the vector spaces $\Tor_i^P(A,\kk)_s$ and $\Tor_i^Q(\kk[S],\kk)_s$ have the same dimension. In particular $\dim_\kk(J/\fm J)_s=\kappa(s)-1$ for $s\in S$, the ideal $J$ is minimally generated by $\eta(S)$ elements, and $(J/\fm J)_s\neq0$ if and only if $s\in\Betti(S)$.
\item The classes of $u^w$, $w\in\Apd$, form a $\kk$-basis of $A$, and $u^w$ is homogeneous of degree $w$. The image of $x^v$ in $A$ is $u^{\wt(v)}$ if $\wt(v)\in\Apd$, and $0$ otherwise.
\item The ideal $J$ is spanned by the monomials $x^v$ with $\wt(v)\notin\Apd$ and the binomials $x^v-x^{v'}$ with $\wt(v)=\wt(v')\in\Apd$. If $\wt(v)\notin\Apd$, then $x^vx^{v'}\in J$ for every $v'\in\N^3$.
\item The socle $\Soc(A)=0:_A\fm$ has the basis $\{u^w : w\in\Apd,\ w+d_i\notin\Apd \text{ for } i=1,2,3\}$, and $\dim_\kk\Soc(A)=t(S)$.
\end{enumerate}
\end{lemma}

\begin{proof}
(a) Since $\kk[S]$ is a domain, $x_0$ is a nonzerodivisor on $\kk[S]=Q/I_S$, and $\kk[S]/x_0\kk[S]=A$, while $Q/x_0Q=P$. If $F_\bullet$ is a minimal graded free resolution of $\kk[S]$ over $Q$, then $H_i(F_\bullet/x_0F_\bullet)=\Tor_i^Q(\kk[S],Q/x_0Q)=0$ for $i>0$, so $F_\bullet/x_0F_\bullet$ is a graded free resolution of $A$ over $P$ with the same shifts. It is minimal, because the entries of its differentials are the images of entries in $\fn$, which lie in $\fm$. Comparing these resolutions gives the first statement. Since $\Tor_1^P(A,\kk)\cong J/\fm J$ and $\Tor_1^Q(\kk[S],\kk)\cong I_S/\fn I_S$, the rest follows from Proposition~\ref{prop:herzog-rosales} and the graded Nakayama lemma.

(b) The ring $\kk[S]$ has the basis $\{u^s : s\in S\}$ and the ideal $u^{d_0}\kk[S]$ has the basis $\{u^s : s\in d_0+S\}$. Hence $A$ has the basis $\{u^s : s\in S\setminus(d_0+S)\}=\{u^w : w\in\Apd\}$, and the image of $x^v$ is the class of $u^{\wt(v)}$.

(c) By (b), each graded component $A_s$ is zero if $s\notin\Apd$ and is spanned by $u^s$ if $s\in\Apd$. The component $P_s$ is spanned by the monomials of weight $s$, and the kernel of $P_s\to A_s$ is therefore $P_s$ in the first case and is spanned by differences of monomials in the second. If $\wt(v)-d_0\in S$, then $\wt(v)+\wt(v')-d_0\in S$.

(d) Since $A$ is graded with graded components of dimension at most one, the graded submodule $\Soc(A)$ is spanned by the elements $u^w$ it contains. By (b), $u^w\in\Soc(A)$ if and only if $w+d_i\notin\Apd$ for $i=1,2,3$. By \eqref{eq:PF}, these $w$ are exactly the elements $f+d_0$ with $f\in\PF(S)$.
\end{proof}

We shall also use the following consequence of the graded Nakayama lemma.

\begin{lemma}\label{lem:nakayama}
Let $G'$ be a finite set of homogeneous generators of $J$. There is a subset $G''\subseteq G'$ that generates $J$ and contains exactly $\dim_\kk(J/\fm J)_s$ elements of degree $s$ for every $s$. In particular $|G''|=\eta(S)$, every $f\in G'\setminus G''$ lies in the ideal generated by $G'\setminus\{f\}$, and if $s\in\Betti(S)$ and $G'$ contains exactly one element of degree $s$, then this element belongs to $G''$.
\end{lemma}

\begin{proof}
The images of the elements of $G'$ of degree $s$ span $(J/\fm J)_s$; choose among them elements whose images form a basis, for each $s$. The resulting set $G''$ generates $J$ by the graded Nakayama lemma, and $|G''|=\dim_\kk J/\fm J=\eta(S)$ by Lemma~\ref{lem:reduction}(a). The last assertion holds because $(J/\fm J)_s\neq0$ for $s\in\Betti(S)$.
\end{proof}

\subsection{Artinian monomial ideals in three variables}

Let $I\subsetneq P$ be a monomial ideal with $\dim_\kk P/I<\infty$. We call such an ideal an \emph{Artinian monomial ideal}. Its \emph{standard set} $\Std(I)=\{v\in\N^3 : x^v\notin I\}$ is a finite nonempty order ideal of $\N^3$; conversely, every finite nonempty order ideal $T$ is the standard set of the Artinian monomial ideal $I_T$ spanned by the monomials $x^v$ with $v\notin T$. A \emph{corner} of $\Std(I)$ is a maximal element. We write $G(I)$ for the set of minimal monomial generators of $I$ and set
\begin{gather*}
\mu(I)=|G(I)|,\qquad t'(I)=\#\{\text{corners of }\Std(I)\},\\
r(I)=\#\{h\in G(I) : xyz \text{ divides } h\}.
\end{gather*}
The socle of $P/I$ is spanned by the classes of the monomials $x^v$ with $v$ a corner, and the Koszul complex identifies $\Tor_3^P(P/I,\kk)$ with this socle up to a shift \cite{BH93}. Hence
\begin{equation}\label{eq:beta3}
t'(I)=\dim_\kk\Soc(P/I)=\dim_\kk\Tor_3^P(P/I,\kk).
\end{equation}
Each variable occurs as a pure power in $G(I)$, so $\mu(I)\ge3$.

\begin{proposition}[{\cite[Theorem~3.2]{Dun18}}]\label{prop:dung}
For every Artinian monomial ideal $I\subseteq\kk[x,y,z]$, one has $\mu(I)\le 2t'(I)+1$.
\end{proposition}

Dung states this for monomial ideals with radical $\fm$, which is our setting, and deduces it from an upper bound for the Betti numbers of monomial ideals in terms of cyclic polytopes, applied to an Alexander dual; compare \cite[Theorem~6.3]{BPS98}. The main theorem of \cite{Dun18} requires an additional hypothesis on a Gr\"obner basis, which does not enter \cite[Theorem~3.2]{Dun18}.

The following complementary inequality is a consequence of the generic deformations and the Scarf triangulations of \cite{BPS98}. We include a proof, because we need a deformation that keeps zero exponents zero, so that the count reflects the supports of the generators.

\begin{proposition}\label{prop:scarf}
For every Artinian monomial ideal $I\subseteq\kk[x,y,z]$, one has $t'(I)\le\mu(I)+r(I)-2$.
\end{proposition}

We recall the relevant notions from \cite{BPS98}. Let $h_1,\dots,h_\mu$ be the minimal generators of a monomial ideal $M$, and for $\sigma\subseteq\{1,\dots,\mu\}$ let $h_\sigma=\lcm(h_i : i\in\sigma)$, with $h_\emptyset=1$. The \emph{Scarf complex} of $M$ is the simplicial complex
\[
\Delta_M=\{\sigma\subseteq\{1,\dots,\mu\} : h_\sigma\neq h_\tau \text{ for all } \tau\neq\sigma\}.
\]
The ideal $M$ is \emph{generic} if no variable occurs with the same positive exponent in two distinct minimal generators. If $\Delta$ is a simplicial complex on $\{1,\dots,\mu\}$ whose vertices are labelled by monomials, $F_\Delta$ denotes the complex of free $P$-modules obtained by homogenizing the simplicial chain complex of $\Delta$ with these labels \cite[Construction~2.1]{BPS98}. By \cite[Lemma~2.2]{BPS98}, $F_\Delta$ is a resolution of $P$ modulo the ideal generated by the labels if and only if, for every monomial $m$, the subcomplex $\{\sigma\in\Delta : h_\sigma \text{ divides } m\}$ is empty or acyclic over $\kk$.

\begin{proof}[Proof of Proposition~\textup{\ref{prop:scarf}}]
Let $\mu=\mu(I)$ and $r=r(I)$, and write $G(I)=\{h_1,\dots,h_\mu\}$ with $h_i=x^{v_i}$ and $v_i=(v_{i1},v_{i2},v_{i3})$.

\emph{Step 1: a deformation preserving supports.} Fix an integer $\nu>\mu$ and put
\[
\tilde v_{is}=\begin{cases} \nu v_{is}+i & \text{if } v_{is}>0,\\ 0 & \text{if } v_{is}=0,\end{cases}\qquad 1\le i\le\mu,\ 1\le s\le3 .
\]
Then the following hold for all $i,j$ and $s$:
\begin{enumerate}[(D1)]
\item $\tilde v_{is}>0$ if and only if $v_{is}>0$;
\item if $v_{is}<v_{js}$, then $\tilde v_{is}<\tilde v_{js}$;
\item if $\tilde v_{is}=\tilde v_{js}>0$, then $i=j$.
\end{enumerate}
Indeed, (D1) is clear. For (D2), if $v_{is}=0$ then $\tilde v_{is}=0<\tilde v_{js}$, and if $0<v_{is}<v_{js}$ then $\tilde v_{is}\le \nu v_{is}+\mu<\nu(v_{is}+1)\le \nu v_{js}<\tilde v_{js}$. For (D3), $\nu v_{is}+i=\nu v_{js}+j$ with $1\le i,j\le\mu<\nu$ forces $i=j$. Let $\tilde h_i=x^{\tilde v_i}$ and let $\tilde I$ be the ideal generated by $\tilde h_1,\dots,\tilde h_\mu$. By (D2), $\tilde v_{is}\le\tilde v_{js}$ implies $v_{is}\le v_{js}$. Hence
\begin{equation}\label{eq:divides}
\tilde h_i \text{ divides } \tilde h_\sigma\ \Longrightarrow\ h_i \text{ divides } h_\sigma \qquad (1\le i\le\mu,\ \sigma\subseteq\{1,\dots,\mu\}),
\end{equation}
since $\tilde v_{is}\le\max_{j\in\sigma}\tilde v_{js}$ for each $s$ gives $v_{is}\le\max_{j\in\sigma}v_{js}$. In particular no $\tilde h_i$ divides another, so $\tilde h_1,\dots,\tilde h_\mu$ are the minimal generators of $\tilde I$. Pure powers are mapped to pure powers, so $\tilde I$ is Artinian; by (D3) it is generic; and by (D1), $\tilde h_i$ is divisible by $xyz$ if and only if $h_i$ is.

\emph{Step 2: a resolution of $P/I$.} Let $\Delta=\Delta_{\tilde I}$ be the Scarf complex of $\tilde I$. Every singleton is a face of $\Delta$, so $\Delta$ has $\mu$ vertices. By \cite[Theorem~3.2]{BPS98}, $F_\Delta$ with the labels $\tilde h_i$ is the minimal free resolution of $P/\tilde I$. We now label the vertex $i$ of $\Delta$ by $h_i$ instead. By \cite[Theorem~4.3 and Corollary~4.4]{BPS98}, relabelling the Scarf complex of a generic deformation in this way gives a free resolution of $P/I$, so that the Betti numbers of $P/I$ are bounded by the face numbers of $\Delta$. A deformation as in \cite[Construction~4.1]{BPS98} may also make zero exponents positive, although some, such as the one in \cite[Example~4.5]{BPS98}, keep them zero. Since we need every zero exponent to remain zero, we give the short argument for our deformation, using \eqref{eq:divides}. Let $m$ be a monomial and let $\rho=\{i : h_i \text{ divides } m\}$. For $\sigma\subseteq\{1,\dots,\mu\}$ we have
\[
h_\sigma \text{ divides } m\iff\sigma\subseteq\rho\iff\tilde h_\sigma \text{ divides } \tilde h_\rho .
\]
The first equivalence is clear. In the second, the implication from left to right is clear; conversely, if $\tilde h_\sigma$ divides $\tilde h_\rho$, then \eqref{eq:divides} shows that $h_i$ divides $h_\rho$, hence $m$, for every $i\in\sigma$. Therefore the subcomplex of faces whose original labels divide $m$ coincides with the subcomplex of faces whose deformed labels divide $\tilde h_\rho$, which is empty or acyclic by \cite[Lemma~2.2]{BPS98} applied to the exact complex $F_\Delta$ with the labels $\tilde h_i$. Applying \cite[Lemma~2.2]{BPS98} again, $F_\Delta$ with the labels $h_i$ is a free resolution of $P/I$. Since a minimal free resolution is a direct summand of every free resolution, \eqref{eq:beta3} gives
\[
t'(I)=\dim_\kk\Tor_3^P(P/I,\kk)\le f_2,
\]
where $f_2$ is the number of $2$-dimensional faces of $\Delta$.

\emph{Step 3: counting triangles.} By \cite[Proposition~5.3 and Construction~5.4]{BPS98}, for every sufficiently large real number $\theta$ the complex $\Delta$ is isomorphic to the complex of faces with a strictly positive inner normal vector of the polytope
\[
\Pi=\operatorname{conv}\{q_i : 1\le i\le\mu\}\subset\R^3,\qquad q_i=(\tilde v_{i1}^{\,\theta},\tilde v_{i2}^{\,\theta},\tilde v_{i3}^{\,\theta}),
\]
the vertex $i$ corresponding to $q_i$; here $0^\theta=0$. By \cite[Corollary~5.5]{BPS98}, $\Delta$ is a regular triangulation of a triangle, and its proof identifies these faces with the faces of $\Pi$ visible from the origin, each of which is a simplex.

We realize this triangulation by the standard radial projection. Let $\Sigma=\{u\in\R^3_{\ge0} : u_1+u_2+u_3=1\}$ and $\pi(u)=u/(u_1+u_2+u_3)$ for $u\in\R^3_{\ge0}\setminus\{0\}$. The polytope $\Pi$ lies in the closed positive orthant, does not contain the origin, and contains a point on each positive coordinate axis, namely the point $q_i$ of each pure power $h_i$. Hence every ray from the origin into the closed positive orthant meets $\Pi$, and the first point of $\Pi$ on such a ray lies on a visible face; conversely, every point of a visible face is the first point of $\Pi$ on its ray. Thus $\pi$ maps the union of the visible faces bijectively, and hence homeomorphically, onto $\Sigma$, and it maps each visible face onto a simplex with the corresponding vertices. Radial projection does not change which coordinates of a point vanish. Hence the vertex $i$ lies on the boundary of $\Sigma$ if and only if $q_i$ has a vanishing coordinate, that is, if and only if $\tilde h_i$, equivalently $h_i$, is not divisible by $xyz$. So the boundary of the triangulated triangle has $b=\mu-r$ vertices and $b$ edges.

Let $f_1$ be the number of edges of $\Delta$. Since $|\Delta|$ is a disk, $\mu-f_1+f_2=1$. Every boundary edge lies in exactly one triangle and every other edge in exactly two, so $3f_2=2f_1-b$. Eliminating $f_1$ gives $f_2=2\mu-b-2=\mu+r-2$, and Step~2 completes the proof.
\end{proof}

\begin{remark}
If $I$ is itself generic, one may take $\tilde I=I$ in the proof; then $F_\Delta$ is the minimal free resolution of $P/I$ by \cite[Theorem~3.2]{BPS98}, and $t'(I)=\mu(I)+r(I)-2$. For $I=(x^2,y^2,z^2,xyz)$ one has $\mu=4$, $r=1$ and $t'=3$, so equality can hold with $r=1$. For $I=(x,y,z)^2$ one has $\mu=6$, $r=0$ and $t'=3<4$.
\end{remark}

\section{Chomicz's staircase}\label{sec:staircase}

\subsection{Minimal relations}

For $i\in\{0,1,2,3\}$ let $a_{ii}$ be the least positive integer $q$ such that $qd_i$ lies in the semigroup generated by the $d_j$ with $j\neq i$. For $m\neq i$ let $b_{m,ii}$ be the least positive integer $q$ such that
\[
qd_i=\sum_{j\neq i}c_jd_j\qquad\text{with } c_j\in\N \text{ and } c_m>0 .
\]
Such $q$ exist, since $d_md_i=d_id_m$, and clearly $a_{ii}\le b_{m,ii}$.

\begin{lemma}[{\cite[Lemma~2.1]{Cho26}}]\label{lem:aii}
Let $i\in\{0,1,2,3\}$ and $c\in\varphi^{-1}(a_{ii}d_i)$. Then either $c=a_{ii}\varepsilon_i'$, where $\varepsilon'_i$ is the $i$-th unit vector of $\N^4$, or $c_i=0$. In particular $\kappa(a_{ii}d_i)\ge2$, so $a_{ii}d_i\in\Betti(S)$.
\end{lemma}

The second sentence follows because $a_{ii}\varepsilon'_i$ is $\cR$-related to no other factorization, while $a_{ii}d_i$ has a factorization with $c_i=0$ by definition.

\begin{definition}[{\cite[Section~2.1]{Cho26}}]\label{def:ordered}
The semigroup $S$ is \emph{primary} if there is $m\in\{0,1,2,3\}$ such that $a_{ii}=b_{m,ii}$ for at least two indices $i\neq m$, and \emph{secondary} otherwise. A \emph{primary ordering} of a primary semigroup is a labelling of its generators as $d_0,d_1,d_2,d_3$ such that
\[
a_{11}=b_{0,11}\quad\text{and}\quad a_{22}=b_{0,22}.
\]
By \cite[Lemma~2.2]{Cho26}, the generators of a secondary semigroup can be labelled so that $a_{00}d_0\neq a_{22}d_2$ and
\begin{gather*}
\varphi^{-1}(a_{00}d_0)=\varphi^{-1}(a_{11}d_1)=\{(a_{00},0,0,0),(0,a_{11},0,0)\},\\
\varphi^{-1}(a_{22}d_2)=\varphi^{-1}(a_{33}d_3)=\{(0,0,a_{22},0),(0,0,0,a_{33})\};
\end{gather*}
such a labelling is a \emph{secondary ordering}. We say that $S$ is \emph{ordered} if its generators carry a primary ordering (when $S$ is primary) or a secondary ordering (when $S$ is secondary).
\end{definition}

The results of \cite{Cho26} that we use are proved for an arbitrary primary ordering; in Section~\ref{sec:choice} we use this freedom. For an ordered semigroup we write $b_{ii}=b_{0,ii}$ for $i=1,2,3$. In a primary ordering $b_{11}=a_{11}$ and $b_{22}=a_{22}$. In a secondary ordering $b_{11}=a_{11}$, because $a_{11}d_1=a_{00}d_0$, while $b_{22}>a_{22}$ and $b_{33}>a_{33}$, because $\varphi^{-1}(a_{22}d_2)=\varphi^{-1}(a_{33}d_3)$ contains no factorization with a positive coefficient of $d_0$.

\subsection{The staircase}

From now on $S$ is ordered. Let
\[
R_S=\{v\in\N^3 : \wt(v)\in\Apd\},\qquad T=\{v\in R_S : v_3<a_{33}\}.
\]
Both are order ideals of $\N^3$, since $\Apd$ is closed under taking summands in $S$. In the language of \cite{Cho26}, $R_S$ is the collection $R$ of unit cubes labelled by elements of the Ap\'ery set, and $T$ is the collection $R\cap\{z\le a_{33}\}$ of \cite[Section~3]{Cho26}.

\begin{proposition}[{\cite[Theorem~2.4 and Propositions~2.7 and 2.8]{Cho26}}]\label{prop:staircase}
If $S$ is ordered, the weight map $\wt$ restricts to a bijection from $T$ onto $\Apd$. In particular $|T|=d_0$.
\end{proposition}

We denote by $I=I_T\subseteq P$ the Artinian monomial ideal with standard set $T$, and we write
\[
\mu=\mu(I),\qquad t'=t'(I),\qquad r=r(I).
\]
These numbers depend on the chosen ordering.

\begin{lemma}\label{lem:basic}
Let $S$ be ordered.
\begin{enumerate}[(a)]
\item Let $i\in\{1,2,3\}$ and $q\ge0$. Then $qd_i\in\Apd$ if and only if $q<b_{ii}$. Consequently $R_S\subseteq[0,b_{11})\times[0,b_{22})\times[0,b_{33})$, and $x^{b_{11}}$, $y^{b_{22}}$ and $z^{a_{33}}$ belong to $G(I)$.
\item If $v\notin T$ and $v_3<a_{33}$, then $\wt(v)\notin\Apd$. In particular every $h\in G(I)$ other than $z^{a_{33}}$ has weight outside $\Apd$.
\item The map $v\mapsto\wt(v)-d_0$ is injective on the set of corners of $T$, and its image contains $\PF(S)$. In particular $t(S)\le t'$.
\end{enumerate}
\end{lemma}

\begin{proof}
(a) If $qd_i\notin\Apd$, write $qd_i=d_0+\sum_jc_jd_j$ with $c_j\in\N$. Since $d_0>0$ we have $c_i<q$, and $(q-c_i)d_i=(1+c_0)d_0+\sum_{j\neq0,i}c_jd_j$ shows $q-c_i\ge b_{ii}$, so $q\ge b_{ii}$. Conversely $b_{ii}d_i\notin\Apd$ by definition, hence $qd_i\notin\Apd$ for $q\ge b_{ii}$. If $v\in R_S$ and $v_i\ge b_{ii}$, then $\wt(v)\in b_{ii}d_i+S\subseteq d_0+S$, a contradiction. Thus $q\varepsilon_1\in T$ if and only if $q<b_{11}$, and similarly for $\varepsilon_2$; since $a_{33}\le b_{33}$, we have $q\varepsilon_3\in T$ if and only if $q<a_{33}$. This proves that the three pure powers are minimal generators of $I$.

(b) The first statement is the definition of $T$. A minimal generator $h\neq z^{a_{33}}$ of $I$ is not divisible by $z^{a_{33}}$, so its exponent vector $v$ satisfies $v\notin T$ and $v_3<a_{33}$.

(c) Let $f\in\PF(S)$ and $w=f+d_0$. By \eqref{eq:PF}, $w\in\Apd$ and $w+d_i\notin\Apd$ for $i=1,2,3$. Let $v\in T$ be the element with $\wt(v)=w$ given by Proposition~\ref{prop:staircase}. Since $\wt(v+\varepsilon_i)=w+d_i\notin\Apd$, we have $v+\varepsilon_i\notin T$ for $i=1,2,3$, so $v$ is a corner. Injectivity on corners follows from Proposition~\ref{prop:staircase}.
\end{proof}

Lemma~\ref{lem:basic}(c) is \cite[Lemma~4.3]{Cho26} in the present language. The next lemma compares $J$ with $I$.

\needspace{8\baselineskip}
\begin{lemma}\label{lem:J}
Let $S$ be ordered and let $a=a_{33}$.
\begin{enumerate}[(a)]
\item If $ad_3\notin\Apd$, then $J=I$. Consequently $\eta(S)=\mu$ and $t(S)=t'$.
\item If $ad_3\in\Apd$, there is a unique monomial $p=x^{p_1}y^{p_2}$ with $(p_1,p_2,0)\in T$ and $\wt(p)=ad_3$. The ideal $J$ is generated by the $\mu$ elements
\[
g=z^a-p\qquad\text{and}\qquad h\in G(I)\setminus\{z^a\}.
\]
Consequently $\eta(S)\le\mu$.
\item One has $ad_3\notin\Apd$ if and only if $a_{33}=b_{33}$. In a secondary ordering, $ad_3\in\Apd$ and $p=y^{a_{22}}$. In a primary ordering with $ad_3\in\Apd$, one has $p_1\ge1$ and $p_2\ge1$.
\end{enumerate}
\end{lemma}

\begin{proof}
(a) By Lemma~\ref{lem:basic}(b) and the hypothesis, every minimal generator of $I$ has weight outside $\Apd$, so $I\subseteq J$ by Lemma~\ref{lem:reduction}(c). Since $\dim_\kk P/I=|T|=d_0=\dim_\kk A=\dim_\kk P/J$, we get $I=J$. Then $\eta(S)=\mu(J)=\mu$ by Lemma~\ref{lem:reduction}(a), and $t(S)=\dim_\kk\Soc(P/I)=t'$ by Lemma~\ref{lem:reduction}(d) and \eqref{eq:beta3}.

(b) By Proposition~\ref{prop:staircase} there is a unique $v\in T$ with $\wt(v)=ad_3$. If $v_3>0$, then $(a-v_3)d_3=v_1d_1+v_2d_2$ with $0<a-v_3<a$, contradicting the definition of $a_{33}$. Hence $v=(p_1,p_2,0)$. Let $J'$ be the ideal generated by $g$ and $G(I)\setminus\{z^a\}$. By Lemma~\ref{lem:reduction}(c) and Lemma~\ref{lem:basic}(b), $J'\subseteq J$. We show that every monomial $x^v$ is congruent modulo $J'$ to $0$ or to a monomial $x^{v'}$ with $v'\in T$, by induction on $v_3$. If $v\in T$ there is nothing to prove. Otherwise $x^v\in I$. If some $h\in G(I)\setminus\{z^a\}$ divides $x^v$, then $x^v\in J'$. If not, $z^a$ divides $x^v$, and $x^v=z^ax^{v-a\varepsilon_3}\equiv p\,x^{v-a\varepsilon_3}$ modulo $J'$; the right-hand side has $z$-exponent $v_3-a$, and induction applies. Hence $\dim_\kk P/J'\le|T|=d_0=\dim_\kk P/J$, and $J'=J$. The bound $\eta(S)\le\mu$ follows from Lemma~\ref{lem:reduction}(a).

(c) By Lemma~\ref{lem:basic}(a), $ad_3\notin\Apd$ if and only if $a\ge b_{33}$, that is, $a=b_{33}$. In a secondary ordering $a_{33}<b_{33}$, so $ad_3\in\Apd$; moreover $(0,p_1,p_2,0)\in\varphi^{-1}(a_{33}d_3)$, which forces $(p_1,p_2)=(0,a_{22})$. In a primary ordering, $(p_1,p_2,0)\in T$ and Lemma~\ref{lem:basic}(a) give $p_1<b_{11}=a_{11}$ and $p_2<b_{22}=a_{22}$. If $p_2=0$, then $p_1d_1=ad_3$ with $0<p_1<a_{11}$, contradicting the definition of $a_{11}$; if $p_1=0$, the same argument applies to $a_{22}$.
\end{proof}

\subsection{The boundary of the staircase}

Let $\{i,j,k\}=\{1,2,3\}$. Following \cite[Section~2.2]{Cho26}, an identity
\begin{equation}\label{eq:22}
\lambda_0d_0+\lambda_id_i=\lambda_jd_j+\lambda_kd_k\qquad\text{with } \lambda_0,\lambda_i,\lambda_j,\lambda_k \text{ positive integers}
\end{equation}
is a \emph{$(2,2)$-type relation}, and the point $Y\in\Z^3$ with $Y_i=-\lambda_i$, $Y_j=\lambda_j$, $Y_k=\lambda_k$ is a \emph{$(2,2)$-type point}. For $Y\in\Z^3$ put $Y^+=(\max\{Y_1,0\},\max\{Y_2,0\},\max\{Y_3,0\})$. For a $(2,2)$-type point, $\wt(Y^+)=\lambda_0d_0+\lambda_id_i$ has a factorization with positive coefficients of $d_0$ and $d_i$, and in particular lies outside $\Apd$.

As in \cite[Section~3]{Cho26}, we fix a set $V\subset\Z^3$ with the following properties.
\begin{enumerate}[(V1)]
\item For every $Y\in V$, the integer $Y_1d_1+Y_2d_2+Y_3d_3$ is a nonnegative multiple of $d_0$.
\item The set $V$ is minimal with respect to inclusion among the sets satisfying (V1) and $T=\N^3\setminus\bigcup_{Y\in V}(Y^++\N^3)$.
\item Every $Y\in V$ with $\wt(Y^+)\notin\{a_{00}d_0,\,b_{11}d_1,\,b_{22}d_2,\,a_{33}d_3\}$ is a $(2,2)$-type point.
\end{enumerate}
Such a set exists by \cite[Theorem~2.4]{Cho26}; see the discussion preceding \cite[Remark~3.2]{Cho26}. (The region associated with $Y$ in \cite{Cho26} consists of the unit cubes with corner in $Y^++\N^3$.) By (V2), $Y\mapsto x^{Y^+}$ is a bijection from $V$ onto $G(I)$: the monomials $x^{Y^+}$ generate $I$, and a point $Y$ for which $x^{Y^+}$ is not a minimal generator, or equals $x^{Y'^+}$ for some $Y'\neq Y$ in $V$, could be removed from $V$. Let
\[
V^*=\{Y\in V : Y \text{ is a } (2,2)\text{-type point}\},\qquad \mathcal N=\{Y\in V^* : \wt(Y^+)\notin\Betti(S)\},
\]
and $U_0=\{\wt(Y^+) : Y\in V^*\setminus\mathcal N\}$, as in \cite[Section~4.2]{Cho26}.

\begin{proposition}[{\cite[(86), Lemmas~4.1 and 4.2]{Cho26}}]\label{prop:boundary}
Let $S$ be ordered. Then
\[
V=V^*\sqcup\{(b_{11},-b_{12},-b_{13}),\ (-b_{21},b_{22},-b_{23}),\ (-a_{31},-a_{32},a_{33}),\ (a_{01},a_{02},a_{03})\},
\]
where $(b_{10},0,b_{12},b_{13})\in\varphi^{-1}(b_{11}d_1)$ and $(b_{20},b_{21},0,b_{23})\in\varphi^{-1}(b_{22}d_2)$ have $b_{10},b_{20}>0$, $(a_{30},a_{31},a_{32},0)\in\varphi^{-1}(a_{33}d_3)$, $(0,a_{01},a_{02},a_{03})\in\varphi^{-1}(a_{00}d_0)$, and the four listed points need not be distinct. Moreover, at most one point of $V\setminus V^*$ has two or more positive coordinates. Consequently:
\begin{enumerate}[(a)]
\item $G(I)$ consists of $x^{b_{11}}$, $y^{b_{22}}$, $z^{a_{33}}$, the monomial $h_0=x^{a_{01}}y^{a_{02}}z^{a_{03}}$ of weight $a_{00}d_0$ (which may coincide with one of the pure powers), and the monomials $x^{Y^+}$ with $Y\in V^*$;
\item each $x^{Y^+}$ with $Y\in V^*$ involves exactly two variables;
\item $r\le1$, and if $r=1$, then $h_0$ is divisible by $xyz$, so that $a_{00}d_0$ has a factorization $(0,c_1,c_2,c_3)$ with $c_1,c_2,c_3>0$.
\end{enumerate}
\end{proposition}

Statements (a)--(c) are immediate from the displayed decomposition: the first three listed points have exactly one positive coordinate and give the pure powers of Lemma~\ref{lem:basic}(a), a $(2,2)$-type point has exactly two positive coordinates, and a full-support generator $x^{Y^+}$ comes from a point $Y$ with three positive coordinates, which can only be $(a_{01},a_{02},a_{03})$; compare \cite[Lemma~4.2]{Cho26}.

\begin{proposition}[{\cite[Lemma~4.13]{Cho26}}]\label{prop:injective}
Let $S$ be ordered. The map $Y\mapsto\wt(Y^+)$ is injective on $V^*\setminus\mathcal N$. In particular $|U_0|=|V^*\setminus\mathcal N|$.
\end{proposition}

\begin{proposition}[{\cite[Lemma~3.8]{Cho26}}]\label{prop:N-primary}
Let $S$ be primary with a primary ordering, and let $Y=(\lambda_1,\lambda_2,-\lambda_3)\in\mathcal N$ with $\lambda_0d_0+\lambda_3d_3=\lambda_1d_1+\lambda_2d_2$ as in \eqref{eq:22}. Then $\lambda_0=a_{00}$ and $\lambda_3<a_{33}$.
\end{proposition}

Chomicz states \cite[Lemma~3.8]{Cho26} for points with negative first coordinate and notes that the analogous statements hold for the other two positions; the case of a negative third coordinate is the one used in the proof of \cite[Lemma~4.11]{Cho26}.

\begin{proposition}[{\cite[Lemmas~3.12, 3.14, 3.15 and 4.16]{Cho26}}]\label{prop:N-secondary}
Let $S$ be secondary with a secondary ordering.
\begin{enumerate}[(a)]
\item No point of $\mathcal N$ has a negative first or second coordinate.
\item If $Y=(\lambda_1,\lambda_2,-\lambda_3)\in\mathcal N$ with $\lambda_0d_0+\lambda_3d_3=\lambda_1d_1+\lambda_2d_2$ as in \eqref{eq:22}, then $\lambda_0=a_{00}$ and $\lambda_3<a_{33}$.
\item $\eta(S)\ge|U_0|+2$.
\end{enumerate}
\end{proposition}

\section{The choice of the distinguished generator}\label{sec:choice}

The numbers $\mu$, $t'$ and $r$ depend on which generator of a primary semigroup plays the role of $d_0$. The following lemma shows that $d_0$ can be chosen so that the case $r=1$ occurs only when $J=I$.

\begin{lemma}\label{lem:choice}
Let $S$ be a primary numerical semigroup with embedding dimension four. Then $S$ has a primary ordering for which $r=0$ or $a_{33}d_3\notin\Apd$.
\end{lemma}

\begin{proof}
Write the generators as $d_0,\dots,d_3$ in some fixed order. The integer $a_{ii}$ depends only on the generator $d_i$; when the generators are relabelled, it moves with $d_i$. Let $\Gamma$ be the directed graph on $\{0,1,2,3\}$ with an arc $i\to j$, for $i\neq j$, whenever some factorization $c$ of $a_{ii}d_i$ with $c_i=0$ has $c_j>0$. Then:
\begin{enumerate}[(O1)]
\item $b_{j,ii}=a_{ii}$ if and only if $i\to j$;
\item every vertex has at least one outgoing arc, since $a_{ii}d_i$ has a factorization with $c_i=0$;
\item a vertex $m$ can serve as $d_0$ in a primary ordering if and only if its indegree is at least two (one then labels two of its in-neighbours as $d_1$ and $d_2$).
\end{enumerate}
Let $\Phi$ be the set of vertices $i$ such that $a_{ii}d_i$ has a factorization $c$ with $c_i=0$ and $c_j>0$ for all $j\neq i$. Every $i\in\Phi$ has outdegree three.

Let $m$ be a vertex of indegree at least two, and consider a primary ordering with $d_0=d_m$. If $m\notin\Phi$, then $r=0$ by Proposition~\ref{prop:boundary}(c). If $m$ has indegree three, then also $a_{33}=b_{33}$ by (O1), so $a_{33}d_3\notin\Apd$ by Lemma~\ref{lem:J}(c). It therefore suffices to find a vertex of indegree at least two that either lies outside $\Phi$ or has indegree three. Since $S$ is primary, there is at least one vertex of indegree at least two.

If $\Phi=\emptyset$, any vertex of indegree at least two will do. If $|\Phi|=4$, every vertex has outdegree three, hence indegree three. If $|\Phi|\in\{2,3\}$, a vertex outside $\Phi$ receives arcs from all vertices of $\Phi$, so it has indegree at least two. Finally let $\Phi=\{v\}$. If some vertex $w\neq v$ has indegree at least two, we choose it. Otherwise each of the three vertices $w\neq v$ has indegree exactly one, its only in-neighbour being $v$, since $v$ has outdegree three. Thus there are no arcs between vertices different from $v$, and by (O2) each $w\neq v$ has an arc to $v$. Hence $v$ has indegree three.
\end{proof}

\section{Counting relations}\label{sec:relations}

In this section we show that, when $J\neq I$ and $r=0$, at most one minimal generator of $I$ is lost in passing to $J$.

\begin{proposition}\label{prop:primary}
Let $S$ be primary with a primary ordering such that $r=0$ and $a_{33}d_3\in\Apd$. Then $\eta(S)\ge\mu-1$.
\end{proposition}

\begin{proof}
Let $a=a_{33}$, let $p=x^{p_1}y^{p_2}$ be as in Lemma~\ref{lem:J}(b), so that $p_1,p_2\ge1$ by Lemma~\ref{lem:J}(c), and let $G'=\{g\}\cup(G(I)\setminus\{z^a\})$, where $g=z^a-p$. By Lemma~\ref{lem:J}(b), $G'$ is a set of $\mu$ homogeneous generators of $J$. Choose $G''\subseteq G'$ as in Lemma~\ref{lem:nakayama} and let $E=G'\setminus G''$. Then $\eta(S)=\mu-|E|$, and we show $|E|\le1$.

\emph{Step 1: elements of $E$.} We have $g\notin E$: the ideal generated by $G'\setminus\{g\}$ is the monomial ideal generated by $G(I)\setminus\{z^a\}$, which does not contain $z^a$, and a monomial ideal containing $g$ would contain $z^a$. Let $h\in E$. Then $h\in G(I)\setminus\{z^a\}$ and $h=\sum_{f\in G''}q_ff$ with $q_f\in P$ and $h\notin G''$. The monomial $h$ occurs with nonzero coefficient on the right-hand side, so it is divisible by a term of some $f\in G''$. These terms are the elements of $G(I)\setminus\{h,z^a\}$, the monomial $z^a$ and the monomial $p$. Since $h$ is a minimal generator of $I$, it is not divisible by another minimal generator of $I$. Hence $p$ divides $h$.

\emph{Step 2: candidates.} Call $h\in G(I)$ a \emph{candidate} if $h=x^\alpha y^\beta$ with $\alpha,\beta\ge1$. By Step~1 and $r=0$, every element of $E$ is a candidate. Let $h=x^\alpha y^\beta$ be a candidate. As $h$ is a minimal generator different from $x^{a_{11}}=x^{b_{11}}$ and $y^{a_{22}}=y^{b_{22}}$, we have $\alpha<a_{11}$ and $\beta<a_{22}$. By Proposition~\ref{prop:boundary}, either $h=h_0$, or $h=x^{Y^+}$ for a $(2,2)$-type point $Y=(\alpha,\beta,-\lambda_3)\in V^*$ with $\lambda_0d_0+\lambda_3d_3=\alpha d_1+\beta d_2$ and $\lambda_0,\lambda_3>0$.

Let $s=\alpha d_1+\beta d_2$ be the degree of a candidate $h$. Then $s\notin\Apd$: if $h=h_0$, then $s=a_{00}d_0$, and otherwise $s=\lambda_0d_0+\lambda_3d_3$. The element $g$ has degree $ad_3\in\Apd$, so $\deg g\neq s$. If $s=a_{11}d_1$, then $\beta d_2=(a_{11}-\alpha)d_1$ with $0<\beta<a_{22}$, contradicting the definition of $a_{22}$; similarly $s\neq a_{22}d_2$. Since $z^a\notin G'$, no element of $G'$ of degree $s$ is $g$ or a pure power. By Proposition~\ref{prop:boundary}(a), every element of $G'$ of degree $s$ is therefore $h_0$ or of the form $x^{Y'^+}$ with $Y'\in V^*$.

\emph{Step 3: $h_0\notin E$.} Suppose that $h_0$ is a candidate, and let $s=a_{00}d_0$, which is a Betti element by Lemma~\ref{lem:aii}. If $Y'\in V^*$, then $\wt(Y'^+)=\lambda'_0d_0+\lambda'_id_i$ for a relation \eqref{eq:22} with $\lambda'_0,\lambda'_i>0$, and a factorization of $a_{00}d_0$ with positive coefficients of $d_0$ and $d_i$ would contradict Lemma~\ref{lem:aii}. By Step~2, $h_0$ is the only element of $G'$ of degree $s$, and $h_0\in G''$ by Lemma~\ref{lem:nakayama}.

\emph{Step 4: Betti boundary generators are not in $E$.} Let $Y=(\alpha,\beta,-\lambda_3)\in V^*\setminus\mathcal N$, so that $x^{Y^+}=x^\alpha y^\beta$ is a candidate, and let $s=\wt(Y^+)=\lambda_0d_0+\lambda_3d_3$, which is a Betti element by the definition of $\mathcal N$. The element $h_0$ does not have degree $s$: otherwise $(\lambda_0,0,0,\lambda_3)$ would be a factorization of $a_{00}d_0$ with positive coefficients of $d_0$ and $d_3$, contradicting Lemma~\ref{lem:aii}. If $Y'\in V^*\setminus\{Y\}$, then $\wt(Y'^+)\neq s$, by Proposition~\ref{prop:injective} if $Y'\notin\mathcal N$, and because $\wt(Y'^+)\notin\Betti(S)$ if $Y'\in\mathcal N$. By Step~2, $x^{Y^+}$ is the only element of $G'$ of degree $s$, and $x^{Y^+}\in G''$ by Lemma~\ref{lem:nakayama}.

\emph{Step 5: at most one exceptional generator.} By Steps~2--4, $E\subseteq\{x^{Y^+} : Y\in\mathcal N,\ Y_3<0\}$. Suppose that $Y=(\lambda_1,\lambda_2,-\lambda_3)$ and $Y'=(\lambda'_1,\lambda'_2,-\lambda'_3)$ are distinct points of $\mathcal N$, both with negative third coordinate. By Proposition~\ref{prop:N-primary},
\[
a_{00}d_0+\lambda_3d_3=\lambda_1d_1+\lambda_2d_2,\qquad a_{00}d_0+\lambda'_3d_3=\lambda'_1d_1+\lambda'_2d_2,
\]
with $0<\lambda_3,\lambda'_3<a_{33}$. Moreover $0<\lambda_1,\lambda'_1<a_{11}$ and $0<\lambda_2,\lambda'_2<a_{22}$, since $x^{Y^+}$ and $x^{Y'^+}$ are candidates (Step~2). Subtracting gives
\[
\delta_1d_1+\delta_2d_2+\delta_3d_3=0,\qquad \delta_1=\lambda_1-\lambda'_1,\quad\delta_2=\lambda_2-\lambda'_2,\quad\delta_3=\lambda'_3-\lambda_3 .
\]
If $\delta_1=\delta_2=0$, then $\delta_3=0$ and $Y=Y'$. Hence some $\delta_i$ is nonzero, and since $d_1,d_2,d_3>0$, the $\delta_i$ take both signs. Interchanging $Y$ and $Y'$ if necessary, we may assume that exactly one $\delta_i$ is positive. Then $\delta_id_i=\sum_{j\neq i}(-\delta_j)d_j$ with $-\delta_j\ge0$, so $\delta_i\ge a_{ii}$ by the definition of $a_{ii}$. On the other hand $\delta_i$ is a difference of two integers in $(0,a_{ii})$, so $\delta_i<a_{ii}$. This contradiction shows that $\mathcal N$ has at most one point with negative third coordinate. Hence $|E|\le1$ and $\eta(S)\ge\mu-1$.
\end{proof}

Step~5 is the argument of \cite[Lemma~4.11]{Cho26}, applied to points of $\mathcal N$ with negative third coordinate.

\begin{proposition}\label{prop:secondary}
Let $S$ be secondary with a secondary ordering. Then $r=0$, $\mathcal N=\emptyset$, $\mu=|V^*|+3$, and $\eta(S)\ge\mu-1$.
\end{proposition}

\begin{proof}
Since $\varphi^{-1}(a_{00}d_0)=\{(a_{00},0,0,0),(0,a_{11},0,0)\}$, Proposition~\ref{prop:boundary}(c) gives $r=0$. In the notation of Proposition~\ref{prop:boundary}, the same equality gives $(a_{01},a_{02},a_{03})=(a_{11},0,0)$; the equality $\varphi^{-1}(a_{11}d_1)=\{(a_{00},0,0,0),(0,a_{11},0,0)\}$ gives $b_{11}=a_{11}$ and $(b_{12},b_{13})=(0,0)$; and $\varphi^{-1}(a_{33}d_3)=\{(0,0,a_{22},0),(0,0,0,a_{33})\}$ gives $(a_{31},a_{32})=(0,a_{22})$. Hence $V\setminus V^*=\{(a_{11},0,0),(-b_{21},b_{22},-b_{23}),(0,-a_{22},a_{33})\}$, whose positive parts are the three distinct pure powers $x^{a_{11}}$, $y^{b_{22}}$, $z^{a_{33}}$. Thus $\mu=|V|=|V^*|+3$.

We show $\mathcal N=\emptyset$. By Proposition~\ref{prop:N-secondary}(a), a point of $\mathcal N$ has the form $Y=(\lambda_1,\lambda_2,-\lambda_3)$, and by Proposition~\ref{prop:N-secondary}(b),
\begin{equation}\label{eq:sec1}
a_{00}d_0+\lambda_3d_3=\lambda_1d_1+\lambda_2d_2,\qquad 0<\lambda_3<a_{33}.
\end{equation}
Moreover $0<\lambda_1<a_{11}$, since $x^{Y^+}=x^{\lambda_1}y^{\lambda_2}$ is a minimal generator of $I$ different from $x^{a_{11}}$. Substituting $a_{00}d_0=a_{11}d_1$ in \eqref{eq:sec1} gives
\[
\lambda_2d_2=(a_{11}-\lambda_1)d_1+\lambda_3d_3,
\]
with $a_{11}-\lambda_1>0$, so $\lambda_2\ge a_{22}$ by the definition of $a_{22}$. Substituting $\lambda_2d_2=(\lambda_2-a_{22})d_2+a_{33}d_3$ in \eqref{eq:sec1} gives
\[
a_{00}d_0=\lambda_1d_1+(\lambda_2-a_{22})d_2+(a_{33}-\lambda_3)d_3 ,
\]
a factorization of $a_{00}d_0$ with positive coefficients of $d_1$ and $d_3$. This contradicts $\varphi^{-1}(a_{00}d_0)=\{(a_{00},0,0,0),(0,a_{11},0,0)\}$. Hence $\mathcal N=\emptyset$.

By Proposition~\ref{prop:injective}, $|U_0|=|V^*|$, and Proposition~\ref{prop:N-secondary}(c) gives $\eta(S)\ge|V^*|+2=\mu-1$.
\end{proof}

\begin{remark}
Chomicz shows that $\mathcal N$ has at most one element in the secondary case \cite[Lemma~4.12]{Cho26}; Proposition~\ref{prop:secondary} shows that it is empty.
\end{remark}

\section{Counting pseudo-Frobenius numbers}\label{sec:type}

We now compare $t(S)$ with the number $t'$ of corners of $T$.

\begin{proposition}\label{prop:topslice}
Let $S$ be ordered with $a_{33}d_3\in\Apd$. Then $t'-t(S)\le r+1$.
\end{proposition}

\begin{proof}
Let $a=a_{33}$ and let $p=x^{p_1}y^{p_2}$ be as in Lemma~\ref{lem:J}(b). By Lemma~\ref{lem:J}(c), $p_2\ge1$ in both the primary and the secondary case. Call a corner $v$ of $T$ \emph{spurious} if $\wt(v)+d_i\in\Apd$ for some $i\in\{1,2,3\}$. By \eqref{eq:PF} and Lemma~\ref{lem:basic}(c), the map $v\mapsto\wt(v)-d_0$ is a bijection from the set of corners that are not spurious onto $\PF(S)$. Hence $t'-t(S)$ is the number of spurious corners.

\emph{Step 1: spurious corners lie on the top slice.} Let $v$ be a corner. For $i=1,2$, the point $v+\varepsilon_i$ is not in $T$ and has third coordinate $v_3<a$, so $\wt(v)+d_i\notin\Apd$ by Lemma~\ref{lem:basic}(b). If $v_3<a-1$, the same applies to $v+\varepsilon_3$, and $v$ is not spurious. Let $v_3=a-1$. Since $ad_3=p_1d_1+p_2d_2$,
\[
\wt(v)+d_3=\wt(v_1,v_2,a)=\wt(v_1+p_1,\,v_2+p_2,\,0).
\]
By the definition of $T$, this weight lies in $\Apd$ if and only if $(v_1+p_1,v_2+p_2,0)\in T$. Thus a corner $v$ is spurious if and only if $v_3=a-1$ and $(v_1+p_1,v_2+p_2,0)\in T$.

\emph{Step 2: the top slice.} Let $D=\{(\alpha,\beta)\in\N^2 : (\alpha,\beta,a-1)\in T\}$. If $D$ is empty there are no spurious corners, so assume $D\neq\emptyset$. It is a finite order ideal of $\N^2$. Since $(\alpha,\beta,a)\notin T$ for all $\alpha,\beta$, the corners of $T$ with third coordinate $a-1$ are exactly the points $(\alpha,\beta,a-1)$ with $(\alpha,\beta)$ maximal in $D$. We list all maximal elements of $D$, including those giving corners that are not spurious, as
\[
(\alpha_0,\beta_0),\ (\alpha_1,\beta_1),\ \dots,\ (\alpha_s,\beta_s),\qquad \alpha_0<\alpha_1<\dots<\alpha_s,\quad \beta_0>\beta_1>\dots>\beta_s .
\]

\emph{Step 3: a generator for each gap.} Fix $i\in\{1,\dots,s\}$. The point $(\alpha_{i-1}+1,\beta_i+1)$ is not in $D$: if it were, it would be dominated by a maximal element $(\alpha_j,\beta_j)$ of $D$ with $\alpha_j>\alpha_{i-1}$ and $\beta_j>\beta_i$, that is, with $j\ge i$ and $j<i$. On the other hand $(\alpha_{i-1},\beta_i+1)\le(\alpha_{i-1},\beta_{i-1})$ and $(\alpha_{i-1}+1,\beta_i)\le(\alpha_i,\beta_i)$ lie in $D$. Choose $h_i\in G(I)$ dividing $x^{\alpha_{i-1}+1}y^{\beta_i+1}z^{a-1}$, and write $h_i=x^{\alpha}y^{\beta}z^{\gamma}$. If $\alpha\le\alpha_{i-1}$, then $h_i$ divides $x^{\alpha_{i-1}}y^{\beta_i+1}z^{a-1}$, whose exponent lies in $T$; this is impossible. Similarly $\beta\le\beta_i$ is impossible. Hence
\[
h_i=x^{\alpha_{i-1}+1}y^{\beta_i+1}z^{\gamma_i},\qquad 0\le\gamma_i\le a-1 .
\]
Different indices $i$ give different $x$-exponents, so $h_1,\dots,h_s$ are distinct.

\emph{Step 4: spurious corners force full support.} Let $i\in\{1,\dots,s\}$ and suppose that the corner $(\alpha_i,\beta_i,a-1)$ is spurious, so that $(\alpha_i+p_1,\beta_i+p_2,0)\in T$ by Step~1. Since $p_1\ge0$, $p_2\ge1$ and $\alpha_i\ge\alpha_{i-1}+1$, this point dominates $(\alpha_{i-1}+1,\beta_i+1,0)$, which therefore lies in $T$. If $\gamma_i=0$, then $h_i=x^{\alpha_{i-1}+1}y^{\beta_i+1}$ would be a monomial of $I$ with exponent in $T$, which is impossible. Hence $\gamma_i\ge1$, and $h_i$ is divisible by $xyz$.

By Steps~3 and~4, the spurious corners $(\alpha_i,\beta_i,a-1)$ with $i\ge1$ correspond injectively to minimal generators of $I$ divisible by $xyz$. Only the corner with $i=0$ is not covered by this injection. Hence there are at most $r+1$ spurious corners.
\end{proof}

\begin{remark}
In the secondary case $r=0$, and Proposition~\ref{prop:topslice} gives $t'\le t(S)+1$, which is \cite[Lemma~4.5]{Cho26}. The assignment $i\mapsto h_i$ of Step~3 is defined for every maximal element of $D$ except the leftmost one $(\alpha_0,\beta_0)$, whether or not the corresponding corner is spurious; Step~4 restricts it to the spurious corners and shows that it then takes values in the generators of full support. Two spurious corners need not be adjacent in the list of Step~2.
\end{remark}

\section{Proof of the main results}\label{sec:proof}

\begin{proof}[Proof of Theorem~\textup{\ref{thm:main}}]
Let $S$ have embedding dimension four. We fix an ordering as follows: if $S$ is secondary, we take a secondary ordering; if $S$ is primary, we take a primary ordering given by Lemma~\ref{lem:choice}. Let $T$, $I$, $\mu$, $t'$ and $r$ be as in Section~\ref{sec:staircase}. Recall that $t(S)\le t'$ (Lemma~\ref{lem:basic}(c)), $t'\le\mu+r-2$ (Proposition~\ref{prop:scarf}) and $\mu\le2t'+1$ (Proposition~\ref{prop:dung}).

\emph{Case 1: $a_{33}d_3\notin\Apd$.} Then $S$ is primary by Lemma~\ref{lem:J}(c), and $\eta(S)=\mu$ and $t(S)=t'$ by Lemma~\ref{lem:J}(a). Since $r\le1$ by Proposition~\ref{prop:boundary}(c),
\[
t(S)=t'\le\mu+r-2\le\mu-1=\eta(S)-1,\qquad \eta(S)=\mu\le 2t'+1=2t(S)+1 .
\]

\emph{Case 2: $a_{33}d_3\in\Apd$.} Then $r=0$: by Proposition~\ref{prop:secondary} if $S$ is secondary, and by the choice of Lemma~\ref{lem:choice} if $S$ is primary. Propositions~\ref{prop:primary} and~\ref{prop:secondary} give $\eta(S)\ge\mu-1$, and Lemma~\ref{lem:J}(b) gives $\eta(S)\le\mu$. Proposition~\ref{prop:topslice} gives $t'\le t(S)+1$. Hence
\[
t(S)\le t'\le\mu-2\le\eta(S)-1,\qquad \eta(S)\le\mu\le2t'+1\le2t(S)+3 .
\]
In both cases $t(S)+1\le\eta(S)\le2t(S)+3$.
\end{proof}

The proof shows slightly more: if $a_{33}d_3\notin\Apd$ for the chosen ordering, then $\eta(S)\le2t(S)+1$.

\begin{proof}[Proof of Corollary~\textup{\ref{cor:betti}}]
By Lemma~\ref{lem:reduction}(a), the total Betti numbers of $\kk[S]$ over $Q$ coincide with those of $A$ over $P$. Since $A$ is Artinian, $\beta_3=\dim_\kk\Tor^P_3(A,\kk)=\dim_\kk\Soc(A)=t(S)$, by the Koszul complex and Lemma~\ref{lem:reduction}(d), and $\beta_i=0$ for $i>3$. Moreover $\beta_0=1$ and $\beta_1=\eta(S)$. The alternating sum of the ranks of a finite free resolution of $\kk[S]$ over $Q$ equals the rank of $\kk[S]$ as a $Q$-module, which is zero; hence $1-\beta_1+\beta_2-\beta_3=0$ and $\beta_2=\eta(S)+t(S)-1$. The inequalities follow from Theorem~\ref{thm:main}.
\end{proof}

\section{Examples and computations}\label{sec:examples}

\begin{example}\label{ex:5678}
Let $S=\langle5,6,7,8\rangle$. Its gaps are $1,2,3,4,9$, so its Frobenius number is $F(S)=9$, and $\PF(S)=\{9\}$; thus $t(S)=1$. The Betti elements are $12,13,14,15,16$, each with exactly two $\cR$-classes; for instance $\varphi^{-1}(12)$ consists of the factorizations $6+6$ and $5+7$, with disjoint supports. Hence $\eta(S)=5=2t(S)+3$.

We follow the proof of Theorem~\ref{thm:main}. With the primary ordering $(d_0,d_1,d_2,d_3)=(5,6,8,7)$ we have $a_{33}=2$, since $14=6+8$, and $\Apd=\{0,6,7,8,14\}$. The staircase is
\[
T=\{(0,0,0),(1,0,0),(0,1,0),(0,0,1),(1,1,0)\},
\]
with corners $(1,1,0)$ and $(0,0,1)$, so $t'=2$, and $G(I)=\{x^2,y^2,z^2,xz,yz\}$, so $\mu=5$ and $r=0$. Here $a_{33}d_3=14\in\Apd$ and $p=xy$, and Lemma~\ref{lem:J}(b) gives
\[
J=(z^2-xy,\ x^2,\ y^2,\ xz,\ yz),
\]
whose generators have degrees $14,12,16,13,15$. The corner $(0,0,1)$ lies on the top slice and is spurious, because $(0+1,0+1,0)\in T$; so $t'-t(S)=1$. Every inequality in the chain $\eta(S)\le\mu\le 2t'+1\le 2t(S)+3$ is an equality.
\end{example}

\begin{remark}\label{rem:ordering}
The choice of Lemma~\ref{lem:choice} matters for the argument, although not for the theorem. For $S=\langle8,9,15,20\rangle$ with the primary ordering $(d_0,d_1,d_2,d_3)=(20,9,15,8)$, one finds $r=1$, $\mu=7$, $t'=5$ and $a_{33}d_3\in\Apd$, while $t(S)=3$ and $\eta(S)=6$. Here $t'-t(S)=2=r+1$, so Proposition~\ref{prop:topslice} is sharp. For this ordering, the general chain of estimates $\eta(S)\le\mu\le2t'+1$ and $t'\le t(S)+r+1$ yields only $\eta(S)\le2t(S)+5$.
\end{remark}

\medskip\noindent\textbf{Computations.} The following finite computations are supplementary; the proofs above do not depend on them. We enumerated all $4{,}771$ numerical semigroups minimally generated by four integers in $\{2,\dots,28\}$. For each of them, $\Ap(S,\min_id_i)$ was computed by shortest paths on residue classes, $t(S)$ by \eqref{eq:PF}, and $\eta(S)=\sum_s(\kappa(s)-1)$, the sum running over the nonzero $s\in S$ with $s\le F(S)+2\max_id_i$, by counting for each such $s$ the connected components of the graph on $\{i : s-d_i\in S\}$ in which $i$ and $j$ are adjacent when $s-d_i-d_j\in S$; this number equals $\kappa(s)$, and beyond the stated bound the graph is complete. All $4{,}771$ semigroups satisfy Theorem~\ref{thm:main}. Table~\ref{tab:range} shows the range of $\eta(S)$ for each type.

\begin{table}[ht]
\caption{Minimum and maximum of $\eta(S)$ for each type $t(S)$, over the $4{,}771$ numerical semigroups minimally generated by four integers in $\{2,\dots,28\}$.}\label{tab:range}
\begin{tabular}{lccccccc}
\toprule
$t(S)$ & 1 & 2 & 3 & 4 & 5 & 6 & 7\\
\midrule
$\min\eta(S)$ & 3 & 4 & 4 & 6 & 7 & 8 & 10\\
$\max\eta(S)$ & 5 & 5 & 7 & 8 & 9 & 10 & 10\\
$2t(S)+3$ & 5 & 7 & 9 & 11 & 13 & 15 & 17\\
\bottomrule
\end{tabular}
\end{table}

For each of these semigroups we also computed the ordering of Lemma~\ref{lem:choice}, the staircase $T$ and the numbers $\mu$, $t'$ and $r$, and checked the inequalities of Lemma~\ref{lem:basic}(c), Propositions~\ref{prop:dung} and~\ref{prop:scarf}, Lemma~\ref{lem:J}, and Propositions~\ref{prop:primary}, \ref{prop:secondary} and~\ref{prop:topslice}. For the $10{,}112$ primary or secondary orderings of the semigroups with generators at most $20$, we checked $r\le1$, Propositions~\ref{prop:dung} and~\ref{prop:scarf}, the bound $\eta(S)\le\mu$, Proposition~\ref{prop:topslice}, and $\eta(S)\ge\mu-1$ whenever $r=0$ and $a_{33}d_3\in\Apd$; the example of Remark~\ref{rem:ordering} was found in this way. We further checked the combinatorial statement in the proof of Lemma~\ref{lem:choice} for all $4{,}087$ pairs consisting of a directed graph on four vertices, in which every vertex has an outgoing arc and some vertex has indegree at least two, and a set $\Phi$ of vertices of outdegree three; and the injection of Proposition~\ref{prop:topslice} for all $979$ nonempty order ideals contained in $\{0,1,2\}^3$, with $0\le p_1\le3$ and $1\le p_2\le3$ ($11{,}748$ cases). The supplementary archive contains the script, written in Python with the standard library only, together with instructions for running it.

\section*{Use of AI-assisted tools}

AI-assisted tools were used in the research and preparation of this manuscript.

\begin{thebibliography}{99}

\bibitem{BPS98}
D.~Bayer, I.~Peeva, and B.~Sturmfels, \emph{Monomial resolutions}, Math. Res. Lett. \textbf{5} (1998), no.~1--2, 31--46.

\bibitem{Bre75a}
H.~Bresinsky, \emph{On prime ideals with generic zero $x_i=t^{n_i}$}, Proc. Amer. Math. Soc. \textbf{47} (1975), 329--332.

\bibitem{Bre75b}
H.~Bresinsky, \emph{Symmetric semigroups of integers generated by $4$ elements}, Manuscripta Math. \textbf{17} (1975), 205--219.

\bibitem{Bre88}
H.~Bresinsky, \emph{Binomial generating sets for monomial curves, with applications in $\mathbb{A}^4$}, Rend. Sem. Mat. Univ. Politec. Torino \textbf{46} (1988), no.~3, 353--370.

\bibitem{BH93}
W.~Bruns and J.~Herzog, \emph{Cohen--Macaulay rings}, Cambridge Studies in Advanced Mathematics, vol.~39, Cambridge University Press, Cambridge, 1993.

\bibitem{CMS24}
G.~Caviglia, A.~Moscariello, and A.~Sammartano, \emph{Bounds for syzygies of monomial curves}, Proc. Amer. Math. Soc. \textbf{152} (2024), 3665--3679.

\bibitem{Cho26}
K.~Chomicz, \emph{The type and cardinality of minimal presentations of numerical semigroups with embedding dimension four}, preprint (2026), arXiv:2609.04000v1.

\bibitem{Dun18}
N.~T.~Dung, \emph{On the type and generators of monomial curves}, Turkish J. Math. \textbf{42} (2018), no.~5, 2112--2124.

\bibitem{FGH87}
R.~Fr\"oberg, C.~Gottlieb, and R.~H\"aggkvist, \emph{On numerical semigroups}, Semigroup Forum \textbf{35} (1987), 63--83.

\bibitem{Her70}
J.~Herzog, \emph{Generators and relations of abelian semigroups and semigroup rings}, Manuscripta Math. \textbf{3} (1970), 175--193.

\bibitem{Kun70}
E.~Kunz, \emph{The value-semigroup of a one-dimensional Gorenstein ring}, Proc. Amer. Math. Soc. \textbf{25} (1970), 748--751.

\bibitem{MS25}
A.~Moscariello and A.~Sammartano, \emph{Open problems on relations of numerical semigroups}, Recent Progress in Ring and Factorization Theory, Springer, Cham, 2025, pp.~365--380.

\bibitem{Ros96}
J.~C.~Rosales, \emph{An algorithmic method to compute a minimal relation for any numerical semigroup}, Internat. J. Algebra Comput. \textbf{6} (1996), no.~4, 441--455.

\bibitem{RGS09}
J.~C.~Rosales and P.~A.~Garc\'ia-S\'anchez, \emph{Numerical semigroups}, Developments in Mathematics, vol.~20, Springer, New York, 2009.

\bibitem{Sta18}
D.~I.~Stamate, \emph{Betti numbers for numerical semigroup rings}, Multigraded Algebra and Applications, Springer, Cham, 2018, pp.~133--157.

\end{thebibliography}

\end{document}
